\hypertarget{chapter-19-the-price-of-a-promise}{%
\section{Chapter 19: The Price of a
Promise}\label{chapter-19-the-price-of-a-promise}}

\begin{quote}
\itshape ``A stable coin is a sentence that must stay true while the whole world
argues with it.''\\
\normalfont\hfill---sigil-usds whitepaper, epigraph
\end{quote}

\hypertarget{scene-1-the-coin-that-promised-to-stand-still}{%
\subsection{Scene 1: The Coin That Promised to Stand
Still}\label{scene-1-the-coin-that-promised-to-stand-still}}

The hardest thing to build, it turned out, was not a coin that no one could print,
nor a coin no one could edit. It was a coin that \emph{promised to stand still.}

``People can't live on Bitcoin's mountain or the Sigil's mathematics,'' Sable
explained to a new cohort, the ones young and angry enough to think every problem
had a clean answer. ``A baker can't price bread in a coin that doubles and halves
by Tuesday. They need a coin pinned to something boring---a coin that says
\emph{one of me is worth one loaf-ish, forever.} A stable coin. And the graveyard
of money is full of stable coins that promised to stand still and then, one bad
afternoon, ran.''

Elena had watched several of those afternoons in her old life. A coin would swear
it was backed, swear each unit was redeemable, and the swearing would hold right up
until enough people tried to redeem at once and discovered the backing was a story.
The promise was always the same. The proof was always missing.

``So how does the Sigil make a promise it can keep?'' a student asked.

``The same way it does everything,'' Elena said. ``It doesn't promise. It
\emph{commits}---and lets you check.''

\hypertarget{scene-2-backed-in-the-open}{%
\subsection{Scene 2: Backed in the
Open}\label{scene-2-backed-in-the-open}}

The Sigil's stable coin did not ask anyone to believe it was backed. Every block,
it committed the proof: the collateral that stood behind every unit, recorded in a
root, checkable in ten milliseconds by a baker or a swarm or a stranger on a train.
You did not trust that the coin could stand still. You watched, continuously, the
exact weight holding it in place. And the supply could not quietly drift, because
every unit minted and every unit burned moved through the same chokepoint that
guarded the native coin---conserved, committed, impossible to inflate in the dark.

``It's the twenty-one million all over again,'' Sable said, ``pointed at a peg
instead of a cap. You can hide who holds the stable coin. You can never hide whether
it's actually backed.''

But a stable coin has a weakness the native coin does not, and Elena named it
before the students could. ``It has to know the price of the boring thing,'' she
said. ``Bread. The dollar. Whatever it's pinned to. And that number lives
\emph{outside} the chain. The whole edifice---the commitment, the conservation, the
checkable backing---rests on one imported fact. The price. And an imported fact is a
door.''

\hypertarget{scene-3-the-oracle-that-blinked}{%
\subsection{Scene 3: The Oracle That
Blinked}\label{scene-3-the-oracle-that-blinked}}

They called the thing that fed the chain the world's prices an oracle, and the
attack, when it came, did not touch the coin at all. It touched the oracle.

For ninety seconds, on a quiet Sunday, the price the chain believed was wrong.
Someone had found a way to make a handful of the oracle's sources agree on a lie---
that the boring thing the coin was pinned to had suddenly, briefly, become worth a
great deal less. The coin's machinery, honest and obedient, did exactly what it was
built to do: it acted to defend a peg against a crisis that did not exist,
liquidating, rebalancing, and---if it had run unchecked for an hour instead of
ninety seconds---transferring a fortune to whoever had whispered the lie.

``The coin didn't fail,'' Sable said, furious and admiring at once. ``The coin did
\emph{exactly} what it promised. It defended the peg. They just lied to it about
which way to lean.''

It was the bridge's seam again, the agent's seam again, the updater's seam again,
wearing yet another suit. The Sigil could make its own internal truth unforgeable.
It could not, by itself, make the \emph{world} honest. Every place the chain had to
import a fact---the other chain's history, an agent's training, the clock, the
price---was a place the seam reopened. The chain was an island of verifiable truth
in an ocean of imported claims, and the oracle was its longest, most fragile bridge
to shore.

\hypertarget{scene-4-many-eyes-on-the-world}{%
\subsection{Scene 4: Many Eyes on the
World}\label{scene-4-many-eyes-on-the-world}}

The fix was the only fix the Sigil had ever had, applied now to the world itself:
never trust one eye. Never trust a few eyes that might agree too easily. Demand the
imported fact be attested by many independent sources, weighted by their long
public record of honesty, with outliers and sudden lurches treated not as
emergencies to act on but as claims to verify before believing. An oracle that
could be made to blink was replaced by one that had to be made to lie in chorus, by
many witnesses who shared no master and had each spent years being right---the
University's reputation, pointed outward, at reality.

``You can't make the world honest,'' Elena told the cohort, ``so you do the next
thing. You make it expensive for the world to lie to you in unison, and you never,
ever act on a single voice. Even about bread. Especially about bread.''

The oracle never blinked again that anyone could profit from. The stable coin stood
still---not because it promised to, but because every block it showed its work, and
every block enough strangers checked. A baker priced bread in it. A swarm settled in
it overnight, between its sips of the cold coin and the heavy one. And the peg held,
not as a vow, but as a verified, continuously re-checked fact.

``A promise you can check,'' Sable said, ``isn't really a promise anymore.''

``No,'' Elena agreed, watching the young witnesses learn to doubt the price of
bread correctly. ``It's something better. It's a fact you're allowed to distrust
until it proves itself. Every block. Forever.''

\hypertarget{chapter-19-notes}{%
\subsection{Chapter Notes}\label{chapter-19-notes}}

\begin{itemize}
\tightlist
\item
  \textbf{A stable coin, backed in the open.} The Sigil's stablecoin commits its
  collateral and conserves its supply through the same chokepoint as the native
  coin, so its backing is continuously verifiable rather than merely promised---the
  21M-cap discipline aimed at a peg.
\item
  \textbf{The oracle problem.} A stablecoin must import an off-chain price; that
  imported fact is the soft target. The attack manipulates the oracle, not the
  coin, making the coin's honest peg-defense machinery act on a lie.
\item
  \textbf{The island and the ocean.} The recurring structural insight: a chain can
  make its \emph{internal} truth unforgeable but must import external facts (other
  chains, agent training, the clock, prices), and every import reopens the seam.
\item
  \textbf{Defense: many independent, reputation-weighted sources.} The fix points
  the network's ``never trust one eye, verify before acting'' discipline outward at
  reality---requiring imported facts to be attested in chorus, with sudden lurches
  verified before they're believed.
\end{itemize}
